% Návrh sekce "Výběr a sběr dat" – k projití a případným úpravám
% Nahradit řádky 1310–1368 v bakalarka2.tex
% Přejmenovat \section{Sběr dat} na \section{Výběr a sběr dat}

\section{Výběr a sběr dat}

\subsection{Sbíraná data}

Pro praktickou část byly vybrány akcie z indexu S\&P500. Tento výběr byl zvolen z důvodu likvidity opčních kontraktů, stability společností a dostupnosti dat z veřejných zdrojů. Cílem bylo získat tituly pokrývající různé úrovně volatility a dividendové politiky.

Historická volatilita byla spočtena z denních uzavíracích cen za posledních 252 obchodních dní jako standardní odchylka logaritmických výnosů přepočtená na roční bázi. Akcie byly rozděleny do tří kategorií: nízká (\(\sigma \le 0{,}35\)), střední (\(0{,}35 < \sigma < 0{,}50\)) a vysoká (\(\sigma \ge 0{,}50\)) volatilita. Dividendový status byl určen podle aktuálního dividendového výnosu \(q\); pokud \(q > 0\), je akcie považována za dividendovou.

\begin{table}[h]
\centering
\caption{Vybrané akcie pro americké opce}
\label{tab:selected_stocks}
\begin{tabular}{lcccc}
\toprule
\textbf{Ticker} & \textbf{Společnost} & \(\sigma\) & \(q\) (\%) & \textbf{Počet opcí} \\
\midrule
\multicolumn{5}{l}{\textbf{Nízká volatilita (\(\sigma \le 0{,}35\))}} \\
JNJ  & Johnson \& Johnson & 0{,}185 & 2{,}16 & 3\,428 \\
KO   & Coca-Cola          & 0{,}167 & 2{,}71 & 2\,738 \\
PG   & Procter \& Gamble  & 0{,}186 & 2{,}95 & 2\,870 \\
MSFT & Microsoft          & 0{,}253 & 0{,}85 & 12\,453 \\
AAPL & Apple              & 0{,}291 & 0{,}39 & 10\,020 \\
GOOGL & Alphabet          & 0{,}304 & 0{,}27 & 13\,970 \\
AMZN & Amazon             & 0{,}335 & 0{,}00 & 8\,321 \\
NFLX & Netflix            & 0{,}335 & 0{,}00 & 25\,755 \\
\midrule
\multicolumn{5}{l}{\textbf{Střední volatilita (\(0{,}35 < \sigma < 0{,}50\))}} \\
META & Meta Platforms     & 0{,}377 & 0{,}34 & 20\,250 \\
NVDA & NVIDIA             & 0{,}412 & 0{,}02 & 15\,611 \\
\midrule
\multicolumn{5}{l}{\textbf{Vysoká volatilita (\(\sigma \ge 0{,}50\))}} \\
TSLA & Tesla              & 0{,}562 & 0{,}00 & 20\,143 \\
AMD  & AMD                & 0{,}624 & 0{,}00 & 11\,841 \\
\bottomrule
\end{tabular}
\end{table}

Z tabulky \ref{tab:selected_stocks} je patrné, že akcie s nízkou volatilitou mají zpravidla vyšší dividendový výnos (JNJ, KO, PG), zatímco akcie s vysokou volatilitou dividendy nevyplácejí (TSLA, AMD). Hranice mezi nízkou a střední volatilitou je u AMZN a NFLX velmi těsná (\(\sigma \approx 0{,}335\)), což je nutné zohlednit při interpretaci výsledků.

Evropské opce jsou zastoupeny opcemi na akciové indexy vybranými z těch, které Yahoo Finance nabízí. Indexy nevyplácejí dividendy, tedy \(q = 0\).

\begin{table}[h]
\centering
\caption{Vybrané indexy pro evropské opce}
\label{tab:selected_indices}
\begin{tabular}{lccc}
\toprule
\textbf{Ticker} & \textbf{Index} & \(\sigma\) & \textbf{Počet opcí} \\
\midrule
\^{}RUT & Russell 2000              & 0{,}220 & 17\,060 \\
\^{}VIX & CBOE Volatility Index     & 1{,}313 & 4\,296 \\
\^{}XDA & DAX                       & 0{,}093 & 35 \\
\^{}XDB & DAX                       & 0{,}069 & 47 \\
\^{}XDE & DAX                       & 0{,}074 & 92 \\
\^{}XDN & DAX                       & 0{,}096 & 66 \\
\bottomrule
\end{tabular}
\end{table}

Index \^{}VIX představuje speciální případ – jeho podkladovým aktivem je volatilita, nikoli cena. Tomu odpovídá i jeho historická volatilita \(\sigma \approx 1{,}3\), která je řádově vyšší než u ostatních indexů. Výsledky pro \^{}VIX je třeba interpretovat s ohledem na tuto skutečnost, neboť standardní oceňovací předpoklady (lognormalita, konstantní volatilita) zde nemusí být splněny.

Německé indexové opce (\^{}XDA{}--\^{}XDN) mají v datasetu velmi omezený počet záznamů (35–92 opcí), což odráží nízkou likviditu těchto kontraktů na Yahoo Finance.

\subsection{Postup sběru}

Sběr dat sestával z několika navazujících kroků, které byly provedeny pomocí skriptů v jazyce Python a R.

\subsubsection{Zdrojová data}

Data byla získána pomocí pythonovské knihovny \texttt{yfinance}, která umožňuje přístup k API rozhraní serveru Yahoo Finance. Pro každý ticker byly staženy následující údaje:
\begin{itemize}
  \item aktuální cena akcie \(S\) v den sběru,
  \item kompletní opční řetězec -- všechna dostupná expirační data, strike ceny, typ (call/put), ceny (bid, ask, last), objem a open interest,
  \item historické denní uzavírací ceny za poslední rok (252 obchodních dní) pro výpočet historické volatility,
  \item dividendová data pro určení spojitého dividendového výnosu \(q\).
\end{itemize}

Sběr proběhl ve čtyřech termínech: 26.~ledna, 27.~února, 27.~března a 26.~května 2026. Opakovaná měření umožňují zachytit vývoj tržních podmínek a otestovat citlivost modelů na změnu ceny podkladového aktiva a bezrizikové úrokové míry. Při dotazování na API byl dodržován odstup nejméně 0,5~s mezi požadavky na různé tickery a 1~s mezi jednotlivými expiracemi, aby nedošlo k zablokování ze strany serveru.

\subsubsection{Zpracování dat}

Historická volatilita \(\sigma\) byla pro každou akcii spočtena z denních logaritmických výnosů
\[
r_i = \ln\frac{S_i}{S_{i-1}},
\]
kde \(S_i\) je uzavírací cena v den \(i\). Směrodatná odchylka těchto výnosů byla přepočtena na roční bázi:
\[
\sigma = \text{std}(r_1,\dots,r_n) \cdot \sqrt{252}.
\]

Spojitý dividendový výnos \(q\) byl určen z historických dividend vyplácených za poslední rok. Pro každou dividendu \(d_i\) vyplacenou v čase \(t_i\) byl vypočten okamžitý výnos \(d_i / S\) a tyto výnosy byly zprůměrovány a přepočteny na spojitou roční míru.

Bezriziková úroková míra \(r\) byla získána z výnosů amerických státních dluhopisů (U.S. Treasury), konkrétně z Daily Treasury Par Yield Curve.\footnote{Dostupné z \url{https://home.treasury.gov/resource-center/data-chart-center/interest-rates}} Pro splatnosti odpovídající jednotlivým expiracím opcí byla úroková míra dopočtena lineární interpolací mezi dvěma nejbližšími publikovanými výnosy.

\subsubsection{Sestavení datasetu}

Data prošla dvěma fázemi zpracování – nejprve v jazyce Python, poté v jazyce R.

\bigskip
\noindent\textbf{Fáze 1: Python (stažení a předzpracování)}

Python skripty obstarávají komunikaci s API serveru Yahoo Finance a základní předzpracování dat. Všechny skripty jsou na sobě nezávislé a spouštějí se samostatně:

\begin{enumerate}
  \item \texttt{price\_fetcher.py} stáhne pro všechny tickery historické denní uzavírací ceny a dividendy za poslední rok. Výstupem je soubor \texttt{stock\_prices\_YYYY-MM-DD.csv}.
  \item \texttt{options\_fetcher\_final.py} stáhne kompletní opční řetězce pro všechny tickery. Ke každé opci ukládá strike cenu, typ (call/put), ceny (bid, ask, lastPrice), objem a open interest. Výstupem jsou soubory \texttt{american\_options\_YYYY-MM-DD.csv} a \texttt{european\_options\_YYYY-MM-DD.csv}.
  \item \texttt{rate\_functioner\_complete.py} provádí kompletní předzpracování – načte raw opční data a ke každé opci dopočítá a přiřadí:
    \begin{itemize}
      \item bezrizikovou úrokovou míru \(r\) z výnosů U.S. Treasury (lineární interpolace podle splatnosti),
      \item aktuální cenu akcie \(S\) z historických cen,
      \item historickou volatilitu \(\sigma\) z denních logaritmických výnosů,
      \item spojitý dividendový výnos \(q\) z historických dividend.
    \end{itemize}
    Zároveň přejmenuje sloupce do konzistentní notace (např. \texttt{historical\_volatility} $\to$ \texttt{sigma}, \texttt{current\_stock\_price} $\to$ \texttt{S}). Výstupem jsou soubory \texttt{american\_options\_YYYY-MM-DD\_greeks.csv} a \texttt{european\_options\_YYYY-MM-DD\_greeks.csv}.
\end{enumerate}

Tato fáze se opakuje pro každý termín sběru dat (leden, únor, březen, květen). Výstupem je sada CSV souborů pro každé měření.

\bigskip
\noindent\textbf{Fáze 2: R (sestavení a enrichování datasetu)}

R skripty načtou výstupy z Python fáze a sestaví z nich jednotný dataset:

\begin{enumerate}
  \item \texttt{data\_processing.rmd} načte pro každý termín sběru:
    \begin{itemize}
      \item stock prices (pro výpočet \(S_{\text{no dividend}}\)),
      \item greeks data (obsahuje všechny potřebné parametry i tržní ceny).
    \end{itemize}
  \item Spočte čas do expirace \(T\) jako \(\text{dny mezi datem sběru a expirací} / 365\).
  \item Vytvoří sloupec \texttt{S\_no\_dividend} (cena akcie navýšená o vyplacené dividendy za poslední rok) pro použití v modelech, které dividendy nezohledňují.
  \item Vypočte poměr \(K/S\) (moneyness) pro snazší kategorizaci opcí.
  \item Sjednotí data ze všech čtyř měření do jednoho datasetu a uloží jej jako \texttt{complete\_dataset.csv}.
\end{enumerate}

Výsledný dataset obsahuje 168~996 opčních kontraktů na 12 amerických akcií a 6 evropských indexů. Čas do expirace \(T\) se pohybuje od 0 do 2{,}8~let a strike ceny \(K\) v rozmezí 0{,}50 až 5~000 jednotek.

\bigskip
\noindent\textbf{Poznámka k tvorbě skriptů:} Python skripty pro stažení a předzpracování dat byly vytvořeny za pomoci generativní umělé inteligence a následně ručně zkontrolovány a upraveny. Veškeré navazující výpočty a analýzy v jazyce R jsou napsány výhradně autorem práce. Všechny skripty jsou přiloženy k práci.
